% ─────────────────────────────────────────────────────────────────
\section{Training Setup}
\label{sec:training}
% ─────────────────────────────────────────────────────────────────

All models are trained with PyTorch~1.13 and PyTorch Geometric~2.4 on a
single NVIDIA GPU (CUDA).
Table~\ref{tab:hyperparams} lists the complete hyperparameter configuration
of the best-performing HierGIN run.

\begin{table}[h]
  \centering
  \caption{HierGIN training hyperparameters (best run: HierGIN-8D).}
  \label{tab:hyperparams}
  \begin{tabular}{ll}
    \toprule
    \textbf{Hyperparameter} & \textbf{Value} \\
    \midrule
    \multicolumn{2}{l}{\textit{Architecture}} \\
    \quad Node backbone          & GIN \\
    \quad Patch backbone         & GIN \\
    \quad Hidden dimension $h$   & 512 \\
    \quad Node GIN layers $L_n$  & 3 \\
    \quad Patch GIN layers $L_p$ & 2 \\
    \quad Patch pooling          & Attention (learned) \\
    \quad Global readout         & Max \\
    \quad Dropout                & 0.30 \\
    \quad RWPE dim               & 0 (disabled) \\
    \midrule
    \multicolumn{2}{l}{\textit{Optimisation}} \\
    \quad Optimiser              & Adam \\
    \quad Learning rate          & $5 \times 10^{-4}$ \\
    \quad Weight decay           & $10^{-4}$ \\
    \quad LR schedule            & Cosine annealing \\
    \quad Warmup epochs          & 5 \\
    \quad Batch size             & 64 \\
    \quad Epochs                 & 80 (patience = 30) \\
    \midrule
    \multicolumn{2}{l}{\textit{Loss and sampling}} \\
    \quad Loss function          & Cross-entropy \\
    \quad Class weighting        & Inverse square-root: $w_c = 1/\sqrt{n_c}$ \\
    \quad Balanced sampler       & Disabled \\
    \midrule
    \multicolumn{2}{l}{\textit{Data}} \\
    \quad Graph nodes $N_v$      & 1\,000 (FPS) \\
    \quad k-NN edges             & $k = 6$ \\
    \quad Patches $k_p$          & 64 (Geodesic Voronoi) \\
    \quad Feature mode           & 8D (all features) \\
    \quad Scalar clamping        & $[\!-5, 5\!]$ \\
    \bottomrule
  \end{tabular}
\end{table}

\subsection{Cosine Annealing with Warmup}

The learning rate schedule follows a linear warmup over the first 5 epochs,
followed by cosine decay to zero:
\begin{equation}
  \eta_t =
  \begin{cases}
    \eta_\text{max} \cdot t / t_\text{warm} & t \leq t_\text{warm}, \\
    \frac{\eta_\text{max}}{2}\!\left(1 + \cos\!\left(\frac{\pi\,(t - t_\text{warm})}{T - t_\text{warm}}\right)\right) & t > t_\text{warm},
  \end{cases}
  \label{eq:lr_schedule}
\end{equation}
where $T = 80$ epochs and $t_\text{warm} = 5$.
Cosine annealing outperforms step decay by approximately $+$0.3--0.5\,pp
macro-F1 in ablations (Section~\ref{sec:ablation}).

\subsection{Class Weighting}

With a 1\,027$\times$ class imbalance, unweighted cross-entropy drives the
model to learn the majority class.
Inverse-square-root weighting assigns:
\begin{equation}
  w_c = \frac{1}{\sqrt{n_c}} \cdot \frac{\sum_{c'} 1/\sqrt{n_{c'}}}{|\mathcal{C}|},
  \label{eq:class_weight}
\end{equation}
normalised so that weights average to 1.
This is a softer weighting than the strict inverse-frequency
$1/n_c$ and avoids over-penalising common classes.
Ablations confirm that inverse-square-root weighting outperforms no weighting
and inverse-frequency weighting.

\subsection{Early Stopping}

Training is stopped if validation macro-F1 does not improve for 30 consecutive
epochs.
The best model checkpoint (by validation F1) is used for test evaluation.
All runs trained for the full 80 epochs in the feature ablation study,
suggesting that the patience criterion was not triggered; longer training was
not explored.

\subsection{Validation Split}

An 80/20 stratified split of the training set is used for validation.
Stratification is over class labels, ensuring each class appears in both
splits proportionally.
The test set labels are not observed during training or model selection.
